\subsection{AI-prompting}
Undervejs i projektet har vi i udstrakt grad brugt AI som sparringspartner til konkrete implementeringer, men også udliciteret opgaver hertil. Det har været smertefrit i det store hele. Vi primært benyttet os af Anthropics Claude og OpenAIs ChatGPT, men også X's Grok. Vi har ligedan genereret assets hermed. Implementeringen af Slimey-systemet er komplet AI's ære. Her bruges Hashtables, som vi kender til, men ikke har anvendt i andre projekter, men det giver god mening, da det giver en kompleksitet på O(1) i modsætning til, hvad der ellers kunne være forfærdeligt. Måske fordi Bevy netop ingen game editor har, er alting jo direkte i kode, og kan således frit kopieres og indsættes til diverse chatbots uden det store bøvl.

Det står klart for os, at det har været vigtigt dog, at vi har en nogenlunde forståelse af, hvordan de forskellige systemer spiller sammen, samt hvad koden gør; det har gjort, at vi nemt har kunnet overflytte elementer efter behov. Ligedan har vi bidraget i nogle tilfælde med god smag og simplifikationer; fordi vi har dybere kendskab til resten af kodebasen, end AI'en, der kun har haft indirekte adgang. Det er ingen tvivl om, at den iterative proces foregår markant hurtigere end hvis processen ikke var AI-enhanced.

Det er dog blevet bemærket, at fordi vores version af Bevy er så ny, har Claude været afdateret, eksempelvis til nogle af fødselssystemerne, hvor den insisterede på at bruge Events, der er forældede og erstattede af Messages, men det har vi så bare fikset. Ligedan oprettede den også vores tilfældighedsgenerator forkert gen$\_$range i stedet for random$\_$range, hvilket selvfølgelig er bagateller.

\subsection{Generelt}
Undervejs i optagelsen af diverse playthroughs, da en jo skal vedhæftes projektet, formåede undertegnede at dø og derfor tabe spillet flere gange. Noget af det skyldes RNG, hvillket blev frustrerende; det er helt bestemt noget, der ville skulle undersøges nærmere. Men hellere, at det er testende, kan mestres og tæmmes, end for nemt. Spillets performance er helt i top til trods for enorme mængder skrald, der alle simuleres i lige så høj grad som spilleren, omend det ikke er de mest krævende fysiske simulationer, hvilket skyldes Bevys evne til at administrere store mængder data simultant (datan låses ikke) samtidig med et foregår via en billig Eulers metode og vektorregning. Det er meget tilfredsstillende, at høste plastik samtidig med, at en sejr faktisk bliver tyngende på grund af sværhedsgraden og føles som et \textit{achievement}. Der er ingen tvivl om, at animationer, der kan bortforklare Slimeys pludselige tilblivelse ville være velkomment. Men spillet får faktisk en vis charme, især på grund af lyd- og musikindtryk, der giver en hel del og sælger arkade-stilen. Spilleren kunne i højere grad godt trænge til feedback, eksempelvis kunne trawleren også eksplodere. Det var også oprindeligt planen, at skibet skulle rotere, når man trykkede A- og D-nøglerne (tilsvarende på controller), men denne lidt langsommere og dårligere evne til at kunne flygte havde simpelthen gjort spillet for svært, så det er godt, at det ikke er med i den endelige version. Man kan måske balancere disse.
